\section{方法}

\subsection{总体流程}

本文采用统一的配置驱动流程，对四种分类方法进行同口径比较：RNA-only、Concat、MOFA、Stacking。流水线由同一入口调度，包含以下步骤：

\begin{enumerate}
    \item 数据读取与样本对齐；
    \item 标签构建与稀有类检查；
    \item 训练折内特征筛选与标准化；
    \item 按方法分支进行融合建模；
    \item 测试折独立评估并汇总统计指标。
\end{enumerate}

该流程的核心约束是：所有可学习步骤均在训练折内完成，测试折仅用于独立推断与评估。

\subsection{四种方法定义}

\subsubsection{RNA-only}

RNA-only 使用表达组学作为单模态基线，用于回答“仅靠 RNA 能达到什么水平”。该分支是后续融合收益评估的参照系。

\subsubsection{Concat（早期融合）}

Concat 将 RNA 与甲基化特征在样本对齐后沿特征维拼接，再输入分类器。该方法实现简单、可解释性直接，是多组学融合的第一层对照。

\subsubsection{MOFA（潜表示融合）}

MOFA 先从多视图数据学习低维潜因子，再将潜因子送入分类器。与直接拼接相比，MOFA 能在潜空间中压缩冗余与噪声。

在本文实现中，MOFA 在训练折拟合；测试折通过训练得到的因子载荷进行投影，从而避免测试样本参与训练拟合。

\subsubsection{Stacking（晚期融合）}

Stacking 在 RNA 与甲基化分支分别训练基模型，使用训练折上的 out-of-fold 概率作为元特征，再训练元分类器输出最终预测。该策略通过分支级决策边界融合提升泛化能力。

\subsection{偏倚控制与可复现性}

为保证方法比较公平，本文采用三项约束：

\begin{enumerate}
    \item 同一数据切分协议下比较四方法；
    \item 同一指标体系（Accuracy、Balanced Accuracy、Macro-F1）；
    \item 同一结果输出格式（fold-level 指标、均值、方差与95\% CI）。
\end{enumerate}

由此可将性能差异主要归因于方法机制差异，而非评估口径差异。
